\loesung{

For code, see the Jupyter notebook or the R solution files.

\begin{enumerate}
    \item Apply a LIME implementation to the FIFA dataset. See notebook for concrete results.

    \item Varying the kernel width $\sigma$ has a strong influence on LIME explanations:
    \begin{itemize}
        \item A large $\sigma$ gives high weight to distant points, making the surrogate model approximate the global behavior of $\fh$ rather than the local behavior around $\xv$. This reduces locality.
        \item A small $\sigma$ concentrates weight on very few nearby points, which can make the surrogate model unstable (high variance across runs due to the random sampling).
        \item There is no principled way to choose $\sigma$; it is a hyperparameter that must be tuned or set by the user.
    \end{itemize}

    \item Key similarities and differences between LIME and SHAP:

    \textbf{Similarities:}
    \begin{itemize}
        \item Both are \emph{local} explanation methods: they explain a single prediction $\fh(\xv)$.
        \item Both produce \emph{additive feature attributions}: each feature receives a numerical contribution to the prediction.
        \item Both are \emph{model-agnostic}: they treat $\fh$ as a black box and only require access to its predictions.
    \end{itemize}

    \textbf{Differences:}
    \begin{itemize}
        \item \textbf{Foundation:} SHAP is grounded in cooperative game theory (Shapley values) and has a unique solution satisfying the efficiency, symmetry, dummy, and additivity axioms. LIME is an optimization-based approach that minimizes a weighted loss with no such axiomatic guarantees.
        \item \textbf{Sampling:} SHAP samples \emph{coalitions} (feature subsets), marginalizing over the remaining features. LIME samples \emph{data points} in the feature space around $\xv$.
        \item \textbf{Surrogate model:} LIME allows flexible surrogate model classes (linear models, decision trees, etc.). SHAP always produces a linear (additive) decomposition.
        \item \textbf{Kernel:} LIME uses a user-chosen exponential kernel with width $\sigma$, making results sensitive to this choice. KernelSHAP uses a specific SHAP kernel that yields exact Shapley values in the limit.
        \item \textbf{Efficiency:} SHAP satisfies the efficiency property: $\sum_j \phi_j = \fh(\xv) - \E[\fh(\Xv)]$. LIME does not guarantee that feature attributions sum to a specific quantity.
        \item \textbf{Stability:} SHAP values are uniquely defined (given the value function), while LIME explanations can vary across runs due to the stochastic sampling of neighborhood points.
    \end{itemize}
\end{enumerate}
}
